\chapter{Algorithmen und Datenstrukturen}

\section{Grundlagen}

\startitemize[packed]
\item Anweisung
\item Sequenz
\item Schleife
\item Verzweigung
\item Struktogramme
\item Effizienz
\startitemize[packed]
\item Speicherbedarf
\item Zahl der Operationen $O(n^2)$
\stopitemize
\item Variablen
\startitemize[packed]
\item Variablentypen
\item Lokal vs global
\item Primitive Datentypen vs Objektreferenzen
\item Anwendungszeck inerhalb Algorithmen
\stopitemize
\item Operationen 
\item Implementieren
\startitemize[packed]
\item {\bf rekursion}
\item Strategie \quotation{Teile und herrsche}
\stopitemize
\stopitemize

\section{Klassen und Objekte}

\startitemize[packed]
\item Klassen
\item Beziehungen
\item Assoziation / Vererbung
\item Klassendiagramme
\stopitemize

\section{Statische vs dynamische Datenstrukturen}

\startitemize[packed]
\item ein- und zweidimenstionale Reihungen
\item Stapel, Schlange und dynamische Reihung
\item Binärbäume
\startitemize[packed]
\item pre-, post- und inorder
\item Suchen udn Einfügen
\stopitemize
\stopitemize

\chapter{Infromationen und Daten}

\section{Kryptologie}

\startitemize[packed]
\item Transposition und Substitution
\item Monoalphabetische Verfahren
\startitemize[packed]
\item Caesar-Verfahren
\stopitemize
\item Häufigkeitsanalyse
\item Beurteilen der Sicherheit
\item Polyalphabetisch: Beispiel Vigenere-Verfahren
\item 
\stopitemize
